%% kapa: exam=1 round=22 year=2011 period=2 form=A subject=appraisal_law subject_ko=부동산관계법규 seq=8 number=48
%% source: Q-Net 감정평가사 기출문제 (cst00309) artlSeq=5036790 file=extracted/1008/5036790/2011년 제22회 감정평가사 1차시험 기출문제/2교시.hwp
%% conversion: kordoc 4.24.51 -> tools/gen1.py (auto); review=proofread-gpt-astra; flags=none
\begin{question}{48}
국토의 계획 및 이용에 관한 법령상 지구단위계획의 지정에 관한 설명으로 옳지 않은 것은?
\begin{choices}
  \item 「도시개발법」에 따라 지정된 도시개발구역의 일부는 제1종 지구단위계획구역으로 지정될 수 있다.
  \item 「택지개발촉진법」에 따라 지정된 택지개발예정지구의 전부는 제1종 지구단위계획구역으로 지정될 수 있다.
  % [원문 그대로] '전부 는'의 띄어쓰기 유지.
  \item 「관광진흥법」에 따라 지정된 관광특구의 전부 는 제1종 지구단위계획구역으로 지정될 수 있다.
  \item 개발제한구역에서 해제되는 구역 중 계획적인 관리가 필요한 지역은 제1종 지구단위계획구역으로 지정될 수 있다.
  \item 계획관리지역으로서 주거개발진흥지구로 지정된 지역은 제1종 지구단위계획구역으로 지정될 수 있다.
\end{choices}
\end{question}
\answer{5}
